\section*{अध्याय \ref{ch:g12:continental-crust} --- महाद्वीपीय भूपर्पटी}

\begin{solution}{exo:g12:continental-crust:1}
\omterm{def:g12:continental-crust:continental-crust}{महाद्वीपीय भूपर्पटी}: 30 से \qty{70}{km}, औसतन लगभग \qty{41}{km};
अवसादी आवरण के नीचे ग्रेनाइटी चट्टानें, नाइस और दूसरी कायांतरित और \omterm{def:g11:igneous-rocks:volcanic-plutonic}{पातालीय
चट्टानें}; औसतन लगभग \qty{61}{\percent} सिलिका; लगभग
\qty{2830}{kg/m^3}; 4.03 अरब वर्ष तक (कण 4.4)। \omterm{def:g11:ocean-ridges:oceanic-lithosphere}{महासागरीय भूपर्पटी}:
लगभग \qty{7}{km}; बेसाल्ट और \omterm{def:g11:ocean-ridges:gabbro}{गैब्रो}; लगभग \qty{50}{\percent} सिलिका; लगभग
\qty{3000}{kg/m^3}; 200 मिलियन वर्ष से कम।
\end{solution}

\begin{solution}{exo:g12:continental-crust:2}
\qty{60.6}{\percent}: \omterm{def:g11:igneous-rocks:felsic-mafic}{मध्यवर्ती}, \omterm{def:g11:igneous-rocks:rock-names}{डायोराइट} की तरह। \qty{66.6}{\percent}:
\omterm{def:g11:igneous-rocks:felsic-mafic}{फ़ेल्सिक}, ग्रेनाइट की तरह।
\end{solution}

\begin{solution}{exo:g12:continental-crust:3}
\omterm{def:g11:ocean-ridges:oceanic-lithosphere}{महासागरीय भूपर्पटी}, लगभग \qty{3000}{kg/m^3}, स्थलमंडलीय प्रावार पर टिकी है,
और पुराना, ठंडा \omterm{def:g11:ocean-ridges:oceanic-lithosphere}{महासागरीय स्थलमंडल} समग्र रूप से डूबने लायक घना है।
\omterm{def:g12:continental-crust:continental-crust}{महाद्वीपीय भूपर्पटी}, लगभग \qty{2830}{kg/m^3}, प्रावार
(लगभग \qty{3380}{kg/m^3}) से कहीं हल्की, बहुत उत्प्लावक है: किसी खाई पर पहुँचने वाला महाद्वीप
\omterm{def:g11:subduction-zones:subduction}{क्षेपण} को रोक देता है।
\end{solution}

\begin{solution}{exo:g12:continental-crust:4}
(1) किसी महाद्वीप का पुराना, स्थिर क्रोड, जो अब विकृत नहीं होता। (2) एक अरब वर्ष से
अधिक। (3) कायांतरित और \omterm{def:g11:igneous-rocks:volcanic-plutonic}{पातालीय चट्टानें}, पूर्व पर्वतों के
गहरे \omterm{def:g12:continental-crust:root}{मूल}।
\end{solution}

\begin{solution}{exo:g12:continental-crust:5}
$4030/190 \approx 21$ गुना। महासागरीय तली कटकों पर बनती है और
\omterm{def:g11:subduction-zones:subduction}{क्षेपण क्षेत्रों} पर नष्ट होती है: उसका कोई भी भाग 200 मिलियन वर्ष से बहुत अधिक
नहीं बचता।
\end{solution}

\begin{solution}{exo:g12:continental-crust:6}
(1) $0.306 \times 210.4 \approx \num{64.4}\text{ मिलियन }\unit{km^2}$।
(2) $210.4 - 64.4 = \num{146.0}\text{ मिलियन }\unit{km^2}$।
(3) $4\pi \times 6371^2 \approx \qty{5.10e8}{km^2}$; $210.4/510.1 \approx
0.41$: \qty{41}{\percent}।
\end{solution}

\begin{solution}{exo:g12:continental-crust:7}
(1) $41100 \times 2830 \approx \qty{1.16e8}{kg}$। (2) $7100 \times 3000 =
\qty{2.13e7}{kg}$। (3) लगभग 5.5 गुना।
\end{solution}

\begin{solution}{exo:g12:continental-crust:8}
\omterm{def:g12:continental-crust:root}{मूल}: $2.5 \times 2830/550 \approx \qty{12.9}{km}$। भूपर्पटी
$2.5 \times 3380/550 \approx \qty{15.4}{km}$ अधिक मोटी है: कुल मिलाकर लगभग \qty{50}{km},
पर्वत-पट्टियों के नीचे पाए गए 50 से \qty{70}{km} के भीतर।
\end{solution}

\begin{solution}{exo:g12:continental-crust:9}
(1) लगभग $0.103 + 0.100 + 0.102 \approx 0.30$: \qty{30}{\percent}।
(2) सबसे नई पट्टी, 0 से 10 मिलियन वर्ष (लगभग \qty{10}{\percent}),
20 से 30 के लगभग बराबर। (3) पुरानी तली \omterm{def:g11:subduction-zones:subduction}{क्षेपण क्षेत्रों} पर नष्ट हो चुकी है;
तली जितनी पुरानी, उसका उतना ही अधिक भाग जा चुका है।
\end{solution}

\begin{solution}{exo:g12:continental-crust:10}
$120 = 0.694\,h - 0.306 \times 100$, इसलिए $h = 150.6/0.694 \approx
\qty{217}{m}$। केवल स्थल को गिनने से महाद्वीपों के निचले, जलमग्न किनारे
छूट जाते हैं, और माध्य कहीं अधिक निकलता है।
\end{solution}

\begin{solution}{exo:g12:continental-crust:11}
द्वीप: लावाओं के नीचे \omterm{def:g11:ocean-ridges:gabbro}{गैब्रो} और \qty{8}{km} पर \omterm{def:g11:seismic-structure:moho}{मोहो}, \omterm{def:g11:ocean-ridges:oceanic-lithosphere}{महासागरीय भूपर्पटी}
(द्वीप की ऊँचाई इसे नहीं बदलती)। मैदान: नाइस और
\qty{38}{km} पर \omterm{def:g11:seismic-structure:moho}{मोहो}, \omterm{def:g12:continental-crust:continental-crust}{महाद्वीपीय भूपर्पटी}।
\end{solution}

\begin{solution}{exo:g12:continental-crust:12}
(1) सटीक: $\tfrac{4}{3}\pi(6371^3 - 6329.9^3) \approx
\qty{2.0829e10}{km^3}$; सन्निकट: $4\pi \times 6371^2 \times 41.1
\approx \qty{2.0964e10}{km^3}$। (2) सन्निकटन लगभग
\qty{0.65}{\percent} अधिक बड़ा है, $H/R \approx 0.0065$ के निकट: पतली परत के लिए
त्रुटि नगण्य है।
\end{solution}

\begin{solution}{exo:g12:continental-crust:13}
(1) $60.6 = (66.6 + 2x)/3$, इसलिए $x = (181.8 - 66.6)/2 = \qty{57.6}{\percent}$।
(2) \omterm{def:g11:igneous-rocks:felsic-mafic}{मध्यवर्ती}। (3) गहरी \omterm{def:g12:continental-crust:continental-crust}{महाद्वीपीय भूपर्पटी} ऊपरी भूपर्पटी से सिलिका में कम
समृद्ध है: अधिक \omterm{def:g11:igneous-rocks:felsic-mafic}{मैफ़िक} खनिज, लोहे, मैग्नीशियम और कैल्सियम से अधिक समृद्ध।
\end{solution}

\begin{solution}{exo:g12:continental-crust:14}
(1) $3 \times 10^8/3 = 10^8$ वर्ष, 100 मिलियन वर्ष। (2) महासागरीय तली की
आयुओं की ही कोटि: तीन-चौथाई 100 मिलियन वर्ष से नई,
कोई भी लगभग 190 से पुरानी नहीं। (3) उसे उतनी ही तेज़ी से नष्ट होना चाहिए जितनी तेज़ी से वह बनती है:
वह \omterm{def:g11:subduction-zones:subduction}{क्षेपण क्षेत्रों} पर वापस प्रावार में डूब जाती है।
\end{solution}

\begin{solution}{exo:g12:continental-crust:15}
(1) $8.6 \times 10^9 / 4.0 \times 10^9 \approx \qty{2.2}{km^3}$ प्रति वर्ष।
(2) स्थिर दर केवल एक संभावना है: भूपर्पटी जल्दी तेज़ी से बढ़ी हो सकती है
और तब से पुनर्चक्रित हुई हो, या देर से बढ़ी हो। सबसे पुरानी चट्टानें और कण सिद्ध करते हैं
कि कुछ भूपर्पटी जल्दी मौजूद थी; वे यह नहीं बता सकते कि कितनी, क्योंकि पुरानी भूपर्पटी
अपरदित हो सकती है या वापस प्रावार में ले जाई जा सकती है और कम चिह्न छोड़ सकती है।
\end{solution}

\begin{solution}{pb:g12:continental-crust:1}
\textbf{1.} $4\pi \times 6371^2 \approx \qty{5.10e8}{km^2}$।

\textbf{2.} $210.4/510.1 \approx 0.41$: \qty{41}{\percent}।

\textbf{3.} $0.306 \times 210.4 \approx \num{64.4}\text{ मिलियन }\unit{km^2}$।

\textbf{4.} $510.1 - 210.4 \approx \num{299.7}\text{ मिलियन }\unit{km^2}$,
\qty{59}{\percent}।

\textbf{5.} महाद्वीपों के किनारे उथले समुद्रों के नीचे,
महाद्वीपीय मग्नतटों पर, हैं।

\textbf{6.} मोटाई, \qty{41}{km}, त्रिज्या, \qty{6371}{km}, की तुलना में
बहुत छोटी है: वक्रता आयतन को एक प्रतिशत से भी
कम बदलती है।

\textbf{7.} $210.4 \times 10^6 \times 41.1 \approx \qty{8.65e9}{km^3}$।

\textbf{8.} \qty{8.65e18}{m^3}।

\textbf{9.} $210.4 \times 10^6 \times 36 \approx \qty{7.57e9}{km^3}$।

\textbf{10.} $(8.65 - 7.57)/7.57 \approx 0.14$: \qty{14}{\percent}।

\textbf{11.} $8.65 \times 10^{18} \times 2830 \approx \qty{2.45e22}{kg}$।

\textbf{12.} $41100 \times 2830 \approx \qty{1.16e8}{kg}$।

\textbf{13.} $299.7 \times 10^6 \times 7.1 \approx \qty{2.13e9}{km^3}$, यानी
\qty{2.13e18}{m^3}; गुणा 3000: लगभग \qty{6.4e21}{kg}।

\textbf{14.} \omterm{def:g12:continental-crust:continental-crust}{महाद्वीपीय भूपर्पटी}, लगभग 3.8 गुना भारी: वह लगभग छह
गुना अधिक मोटी है, जो उसके कम क्षेत्रफल और घनत्व पर भारी पड़ता है।

\textbf{15.} वह ग्रेनाइटी चट्टानों से बनी है, जो क्वार्ट्ज़ और फ़ेल्डस्पारों से समृद्ध हैं, सबसे
हल्के शैलकारी खनिज; \omterm{def:g11:ocean-ridges:oceanic-lithosphere}{महासागरीय भूपर्पटी} बेसाल्ट और
\omterm{def:g11:ocean-ridges:gabbro}{गैब्रो} से बनी है, जो पाइरॉक्सीन और ओलिवीन से समृद्ध हैं।

\textbf{16.} $2.45 \times 10^{22}/5.97 \times 10^{24} \approx 0.0041$।

\textbf{17.} $3 \times 0.41 \times (41.1/6371) \times (2830/5513) \approx
0.0041$: वही।

\textbf{18.} $6.4 \times 10^{21}/5.97 \times 10^{24} \approx 0.0011$: एक प्रतिशत का
लगभग दसवाँ भाग।

\textbf{19.} लगभग \qty{0.4}{\percent}।

\textbf{20.} \omterm{def:g12:continental-crust:continental-crust}{महाद्वीपीय भूपर्पटी} \textbf{पृथ्वी के द्रव्यमान का लगभग \qty{0.4}{\percent}}
है।
\end{solution}
